%%%%%%%%%%%%%%%%%%%%%%%%%%%%%%%%%%%%%%%%%%%%%%%%%%
%%%%		~~~~ Abstrakt ~~~~
%%%%%%%%%%%%%%%%%%%%%%%%%%%%%%%%%%%%%%%%%%%%%%%%%%
% Sollte wirklich kurz sein, 300 Wörter, maximal 500
\begin{abstract}
\justifying % Setzt den Text im Blocksatz
\noindent



Das Testen von Software, insbesondere von regressionskritischen Nutzungsszenarien, ist im Entwicklungsalltag häufig mit einem hohen manuellen und zeitintensiven Aufwand verbunden. Diese Fallstudie untersucht im Kontext der internen Webanwendung \glqq Puma\grqq\ der JUMO GmbH \& Co.\ KG, wie dieser Aufwand durch den Einsatz automatisierter End-to-End-Test-Frameworks zielgerichtet reduziert werden kann. Ziel der Arbeit ist es, auf Basis einer strukturierten Evaluierung das am besten geeignete Werkzeug für die spezifischen, stark reglementierten Unternehmensanforderungen zu identifizieren. Hierfür werden die etablierten Frameworks Selenium und Cypress einander gegenübergestellt.

Die methodische Untersuchung kombiniert eine prototypische Implementierung dreier repräsentativer Testszenarien mit einer anschließenden Nutzwertanalyse nach Kühnapfel. Die Bewertung erfolgt anhand von fünf betrieblich gewichteten Kriterien: Integration in die Infrastruktur (25\,\%), Lernkurve und Syntax (20\,\%), Test-Ausführung und Debugging (20\,\%), Stabilität (20\,\%) sowie Dokumentation und Ökosystem (15\,\%).

Die Ergebnisse der Analyse zeigen, dass Cypress mit einem Gesamtnutzwert von 7,800 aus rein entwicklungstechnischer Sicht dem Framework Selenium (7,200) überlegen ist. Cypress punktet maßgeblich durch eine intuitivere Syntax und signifikant höhere Teststabilität dank integriertem Auto-Waiting. Dennoch offenbart die prototypische Umsetzung gravierende infrastrukturelle Hürden für Cypress: Strenge Proxy-Vorgaben und das etablierte Java-Ökosystem des Unternehmens erschweren eine nahtlose Einbindung erheblich. Selenium lässt sich hingegen als Java-Bibliothek sofort und reibungslos in den bestehenden Maven-Buildprozess integrieren.

Aus diesen Erkenntnissen wird die Handlungsempfehlung abgeleitet, kurzfristig Selenium für die Automatisierung der Tests zu nutzen, um zeitnah erste Automatisierungsgewinne zu erzielen. Mittelfristig sollte das Unternehmen jedoch prüfen, ob die CI/CD-Infrastruktur für Node.js-Umgebungen geöffnet werden kann, um zukünftig von der technischen Überlegenheit von Cypress zu profitieren.

\end{abstract}